% frontmatter/mapa-del-curso.tex -- el año entero en una tabla (el plan
% completo, semana a semana, está en notes/01-plan.md).
\section*{Mapa del curso}

\begin{center}
\renewcommand{\arraystretch}{1.4}
\begin{tabularx}{\linewidth}{@{}c c c c X@{}}
\toprule
\textbf{Trimestre} & \textbf{Días} & \textbf{Época} & \textbf{Dinero} & \textbf{Las ideas} \\
\midrule
1 & 1--65 & otoño & hasta 20 € & El valor de las cosas, el trueque y el
dinero; cuando no hay para todos; lo que es de todos; necesitar y
querer; ahorrar; el precio. \\
2 & 66--130 & invierno & hasta 50 € & Elegir; el dinero de antes; quien
produce y quien compra; cuidar las cosas; el presupuesto de una fiesta;
hacerlo o comprarlo; los servicios; ahorrar para una meta. \\
3 & 131--195 & primavera & hasta 100 € & Los precios de antes; lo que
cuesta una receta; arreglar o comprar nuevo; los oficios y el sueldo;
las tiendas del barrio; la hucha; de la flor al tarro; el banco; el
mercadillo. \\
4 & 196--260 & verano & hasta 100 € & Lo importante, primero; comparar
precios; medir y pesar; apuntar las cuentas; el precio de la miel; la
paga de las fiestas; lo que vale y no cuesta dinero; el puesto de
limonada; planear. \\
\bottomrule
\end{tabularx}
\end{center}

\vspace{6mm}
Cada trimestre tiene 13 semanas de 5 días -- también durante las
vacaciones: el cuaderno sigue siendo una página por día de lunes a
viernes en Navidad y en verano, como \emph{Aprendo a leer}. Cada semana
tiene una idea y una palabra, y el dinero crece con el año: hasta 20 €
en otoño, hasta 50 € en invierno y hasta 100 € en primavera y en verano,
siempre en euros enteros.
\newpage
